% Generated by scripts/materialize_unified_paper.py; do not edit.
\section{Quantum constructive Ramsey search: theorem and contribution}\label{qua:sec:introduction}

The elementary Erd\H{o}s--Szekeres proof of the diagonal Ramsey theorem is
constructive \cite{ErdosSzekeres1935}: choose
a pivot, retain its larger colour neighbourhood, and iterate.  On a graph
with $M$ vertices this uses $O(M)$ edge queries, because every retained set is
represented explicitly.  This part shows that the same proof architecture
has a substantially cheaper quantum implementation when the nested candidate
sets are instead represented by their membership predicates.
We use the elementary $4^{K-1}$ sufficient threshold because it supports this
recursion; it is not the best current extremal bound, which has exponential
base strictly below four \cite{CamposGriffithsMorrisSahasrabudhe2026}.

We work in the promised valid-graph adjacency-oracle model.  One coherent
query returns the colour of a requested unordered vertex pair.  The output is
a set of vertices inducing one colour.  The main theorem is the following.

\begin{theorem}[Implicit-majority quantum Ramsey search]
\label{qua:thm:main}
Let $K\ge2$, $N\ge4^{K-1}$, and $0<\eta<1/2$.  Given coherent edge-oracle
access to a simple undirected graph on $N$ vertices, there is a quantum
algorithm which, with probability at least $1-\eta$, outputs a $K$-clique or
a $K$-vertex independent set using
\[
 O\!\left(2^K K\log\frac K\eta\right)
\]
edge queries in the worst case.  The algorithm verifies every returned
witness and may output $\bot$ on the exceptional event.
\end{theorem}

\begin{corollary}\label{qua:cor:succinct}
For a promised simple graph on $N=2^n$ vertices, one can find a clique or
independent set of order $\lfloor n/2\rfloor+1$ using
\[
 O\!\left(
 \sqrt N\log N\log\frac{\log N}{\eta}
 \right)
\]
edge queries.
\end{corollary}

The first idea is to use capped unknown-solution quantum search to sample
from a set specified only by the $O(K)$ pivot constraints accumulated so far.
There are two ways to choose the next colour class.  The cleaner,
estimation-free method samples one uniform vertex and follows the colour of
its edge to the current pivot.  This chooses a class with probability
proportional to its size.  A new survival lemma shows that this size-biased
recursion completes with probability greater than $1/2$ from
$4^{K-1}$ vertices, giving an
$O(2^K K^2\log K\log(1/\eta))$-query algorithm.

The sharper method uses a scale-aware concentration budget.  Sampling a deep
candidate set costs exponentially more than sampling an early one, so we
estimate early majorities more accurately and permit larger errors later.
The sum of all errors is kept below a fixed constant, whereas the work is
balanced across the recursion levels.  A dyadic schedule removes two
avoidable factors of $K$ from the uniform-accuracy analysis and gives
\cref{qua:thm:main}.  The size-biased argument extends directly to any fixed
number of edge colours; we record the multicolour form as
\cref{qua:cor:multicolour}.

\paragraph{Contributions.}
This part makes four precise contributions.
\begin{itemize}
\item It gives a worst-case
$O(2^K K\log(K/\eta))$ quantum-query algorithm for constructive diagonal
Ramsey search, or $\widetilde O(\sqrt N)$ at the succinct Ramsey scale.
\item It introduces two ways to navigate the implicit recursion: a
size-biased rule with a short survival proof, and a sharper scale-aware
majority estimator whose accuracy budget is matched to the cost of each
level.
\item It extends the size-biased mechanism to a fixed number of colours and
places the result against a randomized classical lower bound of
$\Omega(N^{1-1/\sqrt2})$ in the same parameterization.
\item It gives an explicit parameter audit of the closest published
quantum-query lower-bound claim and isolates an apparent mismatch in the
cited parameter substitution that any comparison in this model must resolve.
\end{itemize}

The result is algorithmic rather than extremal.  It does not change any
numerical Ramsey bound or construct an avoidance colouring.  It gives a
square-root query improvement, up to polylogarithmic factors, over the
standard $O(N)$ constructive recursion, but no matching lower bound against
every randomized classical algorithm is presently known.

There is also a nontrivial literature issue.  Jain, Li, Robere, and Xun state
an $N^{1-o(1)}$ quantum-query lower bound for the $N=2^n$, $K=n/2$
default Ramsey problem in the unnumbered ``Query complexity of RAMSEY''
paragraph of Section III \cite{JainLiRobereXun2024}.  Their hard instances are valid
graph-hash products, while \cref{qua:thm:main} gives an
$O(\sqrt N\operatorname{polylog}N)$ upper bound on promised valid graphs.
As detailed in \cref{qua:sec:prior-art}, direct substitution into the cited
reduction and the fixed-multiplicity Liu--Zhandry theorem yields at most an
$N^{1/24}$ lower bound through that route \cite{LiuZhandry2019}.  We do not
use the reported near-linear consequence.  Our observation concerns this
parameter substitution, not the principal TFNP separation theorems of
Jain et al., and we leave open whether another intended normalization
resolves the discrepancy.

The rest of the paper defines the oracle model and capped sampler, proves the
recurrence and query bound, derives a randomized classical lower bound from
the random-Painter theorem, and records the novelty and nonclaim boundaries.
